% =============================================================================
% 05-disposicao.tex - PARTE 5, ele pode sair daqui? Frames 30 a 35.
% Detalhe operacional na fala e no handout. Na tela, o que decide.
% =============================================================================

% ----------------------------------------------------------------- frame 30 --
\begin{frame}{\ftit{\faQuestion}{O fluxo, comandado por uma pergunta}}
\CRar{6}
\begin{center}
\begin{tikzpicture}[node distance=0.30cm]
  \node[musialBlue, text width=7.6cm, font=\scriptsize, inner ysep=6pt] (q)
    {\sffamily\bfseries \'E poss\'ivel construir seguran\c{c}a suficiente fora de
     um ambiente protegido neste momento?};

  \node[crnode={campoverde}{4.25cm}, font=\scriptsize, inner ysep=4pt,
        below left=0.62cm and -0.2cm of q.south] (sim)
    {\textbf{Sim}\\[1pt] restri\c{c}\~ao de acesso pactuada, plano de
     seguran\c{c}a escrito, rede acionada, seguimento com data e contato ativo};
  \node[crnode={campobordo}{4.25cm}, font=\scriptsize, inner ysep=4pt,
        below right=0.62cm and -0.2cm of q.south] (nao)
    {\textbf{N\~ao}\\[1pt] perman\^encia em ambiente protegido, laudo
     circunstanciado, modalidade definida pelo consentimento, seguran\c{c}a
     ambiental at\'e a vaga};

  \draw[crseta=campoverde] (q.south) -- ++(0,-0.28) -| (sim.north);
  \draw[crseta=campobordo] (q.south) -- ++(0,-0.28) -| (nao.north);

  \coordinate (baixo) at ($(sim.south)!0.5!(nao.south)$);
  \node[crnode={campoambar}{9.2cm}, font=\scriptsize, inner ysep=4pt,
        below=0.55cm of baixo] (re)
    {Em qualquer das duas sa\'idas: reavaliar. A resposta vale para este momento,
     e n\~ao para a semana.};
  \draw[musialThin] (sim.south) |- ($(re.north)+(0,0.26)$) -- (re.north);
  \draw[musialThin] (nao.south) |- ($(re.north)+(0,0.26)$) -- (re.north);
\end{tikzpicture}
\end{center}

\mensagem{N\~ao existe uma terceira sa\'ida chamada orientado e liberado.}
\end{frame}

% ----------------------------------------------------------------- frame 31 --
\begin{frame}{\ftit{\faAmbulance}{Quando remover}}
\CRar{1.6}
\begin{columns}[onlytextwidth, t]
\begin{column}{0.47\textwidth}
\begin{caixaagir}{Remover para servi\c{c}o com retaguarda}
\begin{itemc}
\item risco cl\'inico da tentativa em curso
\item intoxica\c{c}\~ao que a unidade n\~ao sustenta
\item agita\c{c}\~ao ou psicose sem condi\c{c}\~oes de manejo local
\item necessidade de ambiente protegido sem leito na unidade
\end{itemc}
\end{caixaagir}
\end{column}
\calha
\begin{column}{0.47\textwidth}
\begin{caixaproteger}{Antes de a ambul\^ancia chegar}
\begin{itemc}
\item acompanhante definido
\item seguran\c{c}a ambiental no intervalo
\item o que j\'a foi feito, escrito
\item quem recebe do outro lado, avisado
\end{itemc}
\end{caixaproteger}
\end{column}
\end{columns}

\CRvg
\begin{caixaagir}{Enquanto espera}
O intervalo entre a decis\~ao e a remo\c{c}\~ao \'e parte da conduta, e n\~ao
uma pausa administrativa. Paciente em espera de vaga continua sob
observa\c{c}\~ao e continua sendo reavaliado.
\end{caixaagir}
\end{frame}

% ----------------------------------------------------------------- frame 32 --
% Sequencia fixa e LITERAL na tela.
\begin{frame}{\ftit{\faPhoneVolume}{O que informar \`a regula\c{c}\~ao}}
\CRar{6}
\begin{center}
\begin{tikzpicture}[node distance=0.13cm]
  \node[crnode={campoazul}{8.9cm}, font=\scriptsize, inner ysep=2.6pt] (a)
    {\textbf{1.} quem \'e o paciente: nome, idade, de onde est\'a falando};
  \node[crnode={campoazul}{8.9cm}, font=\scriptsize, inner ysep=2.6pt, below=of a] (b)
    {\textbf{2.} o que aconteceu: o motivo do atendimento, em uma frase};
  \node[crnode={campoazul}{8.9cm}, font=\scriptsize, inner ysep=2.6pt, below=of b] (c)
    {\textbf{3.} estado mental atual};
  \node[crnode={campoazul}{8.9cm}, font=\scriptsize, inner ysep=2.6pt, below=of c] (d)
    {\textbf{4.} suicidalidade atual: idea\c{c}\~ao, inten\c{c}\~ao, plano, acesso};
  \node[crnode={campoambar}{8.9cm}, font=\scriptsize, inner ysep=2.6pt, below=of d] (e)
    {\textbf{5.} comportamentos recentes e o que mudou nas \'ultimas horas};
  \node[crnode={campoambar}{8.9cm}, font=\scriptsize, inner ysep=2.6pt, below=of e] (f)
    {\textbf{6.} intoxica\c{c}\~ao, agita\c{c}\~ao ou psicose};
  \node[crnode={campoverde}{8.9cm}, font=\scriptsize, inner ysep=2.6pt, below=of f] (g)
    {\textbf{7.} condutas j\'a realizadas neste atendimento};
  \node[crnode={campobordo}{8.9cm}, font=\scriptsize, inner ysep=2.6pt, below=of g] (h)
    {\textbf{8.} por que o ambiente atual n\~ao \'e suficiente};
\end{tikzpicture}
\end{center}

\mensagem{A oitava \'e a que decide a vaga. Diga-a nesta ordem e n\~ao come\c{c}e
por ela.}
\end{frame}

% ----------------------------------------------------------------- frame 33 --
% Segunda tabela autorizada: as tres modalidades de internacao.
\begin{frame}{\ftit{\faHospital}{As tr\^es modalidades de interna\c{c}\~ao}}
\CRar{0}
\begin{CRtabela}
\scriptsize
\begin{center}
\begin{tabular}{P{2.30cm} P{2.75cm} P{2.90cm} P{2.60cm}}
\toprule
\rowcolor{campografite!12}
& \textbf{Volunt\'aria} & \textbf{Involunt\'aria} & \textbf{Compuls\'oria} \\
\midrule
quem solicita & o pr\'oprio usu\'ario & terceiro & decis\~ao judicial \\
quem autoriza & m\'edico com CRM & m\'edico com CRM & o ju\'iz determina \\
qual documento & termo assinado na admiss\~ao & laudo e formul\'ario &
decis\~ao judicial \\
comunica a quem & n\~ao se aplica & Minist\'erio P\'ublico Estadual & conforme a
decis\~ao \\
em que prazo & n\~ao se aplica & 72 horas, pelo respons\'avel t\'ecnico &
conforme a decis\~ao \\
como termina & a pedido do usu\'ario & a pedido do familiar ou do especialista
respons\'avel & por decis\~ao judicial \\
\bottomrule
\end{tabular}
\end{center}
\end{CRtabela}

\begin{caixanorma}{Vale para as tr\^es}
Toda interna\c{c}\~ao exige \norma{laudo circunstanciado}, a
comunica\c{c}\~ao ao Minist\'erio P\'ublico \norma{se repete na alta}, e a
\agir{compuls\'oria \'e judicial, n\~ao \'e a via da urg\^encia}.
\end{caixanorma}

\fonte{Lei 10.216/2001, arts. 6 a 9; Portaria GM\slash MS 2.391/2002;
Resolu\c{c}\~ao CFM 2.057/2013. \selonorma}
\end{frame}

% ----------------------------------------------------------------- frame 34 --
\begin{frame}{\ftit{\faHandPaper}{Conten\c{c}\~ao}}
\CRar{4.6}
\begin{caixanorma}{Medida excepcional de seguran\c{c}a}
A conten\c{c}\~ao \'e admiss\'ivel quando \'e o meio mais adequado para prevenir
dano imediato ou iminente. N\~ao \'e puni\c{c}\~ao, n\~ao \'e rotina de
admiss\~ao e n\~ao substitui vigil\^ancia.
\end{caixanorma}

\CRvg
\begin{center}
\begin{tikzpicture}[node distance=0.24cm]
  \node[crnode={campografite}{2.05cm}, font=\scriptsize, minimum height=0.90cm] (a)
    {indica\c{c}\~ao registrada};
  \node[crnode={campografite}{2.05cm}, font=\scriptsize, minimum height=0.90cm,
        right=of a] (b) {prescri\c{c}\~ao m\'edica};
  \node[crnode={campografite}{2.05cm}, font=\scriptsize, minimum height=0.90cm,
        right=of b] (c) {monitoriza\c{c}\~ao cont\'inua};
  \node[crnode={campografite}{2.05cm}, font=\scriptsize, minimum height=0.90cm,
        right=of c] (d) {registro em prontu\'ario};
  \draw[crseta=campografite] (a.east) -- (b.west);
  \draw[crseta=campografite] (b.east) -- (c.west);
  \draw[crseta=campografite] (c.east) -- (d.west);
\end{tikzpicture}
\end{center}

\CRvg
\begin{caixaagir}{Reavalia\c{c}\~ao}
Conten\c{c}\~ao tem in\'icio, tem revis\~ao em intervalo curto e tem fim
escrito. Conten\c{c}\~ao sem hora de t\'ermino no prontu\'ario \'e
conten\c{c}\~ao indefinida.
\end{caixaagir}

\fonte{Resolu\c{c}\~ao CFM 2.057/2013, art. 16. \selonorma}
\end{frame}

% ----------------------------------------------------------------- frame 35 --
\begin{frame}{\ftit{\faLock}{Sigilo}}
\CRar{2.5}
\begin{caixanorma}{Sigilo n\~ao \'e sil\^encio absoluto}
Diante de risco relevante \`a vida, a comunica\c{c}\~ao a quem pode proteger \'e
leg\'itima. Ela \'e \norma{proporcional \`a necessidade de prote\c{c}\~ao}: diz o
necess\'ario, a quem pode agir, e n\~ao mais do que isso.
\end{caixanorma}

\CRvg
\begin{columns}[onlytextwidth, t]
\begin{column}{0.47\textwidth}
\begin{caixaproteger}{Como se faz}
\begin{itemc}
\item buscar o consentimento primeiro, e insistir nele
\item dizer ao paciente o que ser\'a dito e a quem
\item informar o necess\'ario para proteger, n\~ao o conte\'udo inteiro
\item registrar a justificativa
\end{itemc}
\end{caixaproteger}
\end{column}
\calha
\begin{column}{0.47\textwidth}
\begin{caixanorma}{Deveres espec\'ificos}
\begin{itemc}
\item crian\c{c}a e adolescente: respons\'avel e Conselho Tutelar quando couber
\item pessoa idosa em situa\c{c}\~ao de viol\^encia ou neglig\^encia: prote\c{c}\~ao devida
\item nunca prometer sigilo absoluto antes de saber o que ser\'a contado
\end{itemc}
\end{caixanorma}
\end{column}
\end{columns}
\end{frame}
